\chapter{光的干涉}

\problem{1.1}{波长为 $500\,\mathrm{nm}$ 的绿光投射在间距 $d$ 为 $0.022\,\mathrm{cm}$ 的双缝上，在距离 $r_0$ 为 $180\,\mathrm{cm}$ 处的光屏上形成干涉条纹，求两个亮条纹之间的距离。若改用波长为 $700\,\mathrm{nm}$ 的红光投射到此双缝上，两个亮纹之间的距离又为多少？算出这两种光第 2 级亮纹位置的距离。}
\answerarea

\problem{1.2}{在杨氏实验装置中，光源波长为 $640\,\mathrm{nm}$，两狭缝间距 $d$ 为 $0.4\,\mathrm{mm}$，光屏离狭缝的距离 $r_0$ 为 $50\,\mathrm{cm}$。试求：（1）光屏上第 1 级亮条纹和中央亮条纹之间的距离；（2）若 $P$ 点离中央亮条纹为 $0.1\,\mathrm{mm}$，问两束光在 $P$ 点的相位差是多少？（3）求 $P$ 点的光强度和中央点的强度之比。}
\answerarea
\clearpage

\problem{1.3}{把折射率为 $1.5$ 的玻璃片插入杨氏实验的一束光路中，光屏上原来第 5 级亮条纹所在的位置变为中央亮条纹，试求插入的玻璃片的厚度。已知光波长为 $600\,\mathrm{nm}$。}
\answerarea

\problem{1.4}{波长为 $500\,\mathrm{nm}$ 的单色平行光射在间距 $d$ 为 $0.2\,\mathrm{mm}$ 的双狭缝上，通过其中一个缝的能量为另一个的 $2$ 倍，在离狭缝 $50\,\mathrm{cm}$ 的光屏上形成干涉图样，求干涉条纹间距和条纹的可见度。}
\answerarea
\clearpage

\problem{1.5}{波长为 $700\,\mathrm{nm}$ 的光源与菲涅耳双镜的相交棱之间距离为 $20\,\mathrm{cm}$，棱到光屏间的距离 $L$ 为 $180\,\mathrm{cm}$，若所得干涉条纹中相邻亮条纹的间隔为 $1\,\mathrm{mm}$，求双镜平面之间的夹角 $\theta$。}
\answerarea

\problem{1.6}{在题 1.6 图所示的劳埃德镜实验中，光源 $S$ 到观察屏的距离为 $1.5\,\mathrm{m}$，到劳埃德镜面的垂直距离为 $2\,\mathrm{mm}$。劳埃德镜长 $40\,\mathrm{cm}$，置于光源和屏之间的中央。（1）若光波波长 $\lambda=500\,\mathrm{nm}$，问条纹间距是多少？（2）确定屏上可以看见条纹的区域大小，此区域内共有几条条纹？（提示：产生干涉的区域 $P_1P_2$ 可由图中的几何关系求得。）\fig[0.62]{1_6}}
\answerarea
\clearpage

\problem{1.7}{试求能产生红光（$\lambda=700\,\mathrm{nm}$）的二级反射亮干涉条纹的肥皂膜厚度，已知肥皂膜折射率为 $1.33$，且平行光与法向成 $30^\circ$ 角入射。}
\answerarea

\problem{1.8}{透镜表面通常镀一层如 $\mathrm{MgF_2}$（$n=1.38$）一类的透明物质薄膜，目的是利用干涉来降低玻璃表面的反射。为了使透镜在可见光谱的中心波长（$550\,\mathrm{nm}$）处产生极小的反射，则镀层最小厚度为多少？}
\answerarea
\clearpage

\problem{1.9}{在两块玻璃片之间一边放一条厚纸，另一边相互压紧。玻璃片 $l$ 长 $10\,\mathrm{cm}$，纸厚为 $0.05\,\mathrm{mm}$，从 $60^\circ$ 的反射角进行观察，问在玻璃片单位长度内看到的干涉条纹数目是多少？设单色光源波长为 $500\,\mathrm{nm}$。}
\answerarea

\problem{1.10}{在上题装置中，沿垂直于玻璃片表面的方向看去，看到相邻两条暗纹间距为 $1.4\,\mathrm{mm}$。已知玻璃片长 $17.9\,\mathrm{cm}$，纸厚 $0.036\,\mathrm{mm}$，求光波的波长。}
\answerarea
\clearpage

\problem{1.11}{波长为 $400\sim 760\,\mathrm{nm}$ 的可见光正射在一块厚度为 $1.2\times 10^{-6}\,\mathrm{m}$、折射率为 $1.5$ 的薄玻璃片上，试问从玻璃片反射的光中哪些波长的光最强。}
\answerarea

\problem{1.12}{迈克耳孙干涉仪的反射镜 $\mathrm{M}_2$ 移动 $0.25\,\mathrm{mm}$ 时，看到条纹移过的数目为 $909$ 个，设光为垂直入射，求所用光源的波长。}
\answerarea
\clearpage

\problem{1.13}{迈克耳孙干涉仪平面镜的面积为 $4\times 4\,\mathrm{cm}^2$，观察到该镜上有 $20$ 个条纹，当入射光的波长为 $589\,\mathrm{nm}$ 时，两镜面之间的夹角为多大？}
\answerarea

\problem{1.14}{调节一台迈克耳孙干涉仪，使其用波长为 $500\,\mathrm{nm}$ 的扩展光源照明时会出现同心圆环条纹，若要使圆环中心处相继出现 $1\,000$ 条圆环条纹，则必须将移动一臂多远的距离？若中心是亮的，试计算第一暗环的角半径。（提示：圆环是等倾干涉图样，计算第一暗环角半径时可利用 $\theta\approx\sin\theta$ 及 $\cos\theta\approx 1-\dfrac{\theta^2}{2}$ 的关系。）}
\answerarea
\clearpage

\problem{1.15}{用单色光观察牛顿环，测得 $j$ 级亮环的直径为 $3\,\mathrm{mm}$，在它的第 $j+5$ 级亮环的直径为 $4.6\,\mathrm{mm}$，所用平凸透镜的凸面曲率半径为 $1.03\,\mathrm{m}$，求此单色光的波长。}
\answerarea

\problem{1.16}{在反射光中观察某单色光所形成的牛顿环，其第 2 级亮环与第 3 级亮环间距为 $1\,\mathrm{mm}$，求第 19 和 20 级亮环之间的距离。}
\answerarea
\clearpage

\sproblemfig[0.32]{1.17}{牛顿环可由两个曲率半径很大的平凸透镜之间的空气层产生（题 1.17 图），平凸透镜 A 和 B 的曲率半径分别为 $R_{\mathrm{A}}$ 和 $R_{\mathrm{B}}$，在波长为 $600\,\mathrm{nm}$ 的单色光垂直照射下观察到反射牛顿环第 10 级暗环半径 $r_{\mathrm{AB}}=4\,\mathrm{mm}$。若另有曲率半径为 $R_{\mathrm{C}}$ 的平凸透镜 C（图中未画出），并且 B、C 组合和 A、C 组合产生的第 10 级暗环半径分别为 $r_{\mathrm{BC}}=4.5\,\mathrm{mm}$ 和 $r_{\mathrm{AC}}=5\,\mathrm{mm}$，试计算 $R_{\mathrm{A}}$、$R_{\mathrm{B}}$ 和 $R_{\mathrm{C}}$。}{1_17}
\answerarea

\sproblem{1.18}{菲涅耳双棱镜实验装置尺寸如下：缝到棱镜的距离为 $5\,\mathrm{cm}$，棱镜到屏的距离为 $95\,\mathrm{cm}$，棱镜角为 $\alpha=179^\circ 32'$，构成棱镜玻璃材料的折射率 $n'=1.5$，采用的是单色光。当厚度均匀的肥皂膜横过双棱镜的一半部分放置，该系统中心部分附近的条纹相对原先有 $0.8\,\mathrm{mm}$ 的位移。若肥皂膜的折射率为 $n=1.35$，试计算肥皂膜厚度的最小值为多少？}
\answerarea
\clearpage

\sproblemfig[0.18]{1.19}{将焦距为 $50\,\mathrm{cm}$ 的会聚透镜中央部分 C 切去（见题 1.19 图），余下的 A、B 两部分仍旧粘起来，C 的宽度为 $1\,\mathrm{cm}$。在对称轴线上距透镜 $25\,\mathrm{cm}$ 处置一点光源，发出波长为 $692\,\mathrm{nm}$ 的激光，在对称轴线上透镜的另一侧 $50\,\mathrm{cm}$ 处置一光屏，屏面垂直于轴线。试求：
\begin{enumerate}
  \item 干涉条纹的间距是多少？
  \item 光屏上呈现的干涉图样是怎样的？
\end{enumerate}}{1_19}
\answerarea

\sproblem{1.20}{将焦距为 $5\,\mathrm{cm}$ 的薄凸透镜 L 沿直线方向剖开（见题 1.20 图）分成两部分 A 和 B，并将 A 部分沿主轴右移至 $2.5\,\mathrm{cm}$ 处，这种类型的装置称为梅斯林对切透镜。若将波长为 $632.8\,\mathrm{nm}$ 的点光源 $P$ 置于主轴上离透镜 B 的距离为 $10\,\mathrm{cm}$ 处，试分析：
\begin{enumerate}
  \item 根据中学几何光学的知识分析成像情况如何？
  \item 若在 B 右边 $10.5\,\mathrm{cm}$ 处置一光屏，则在光屏上观察到的干涉图样如何？
\end{enumerate}\fig[0.45]{1_20}}
\answerarea
\clearpage

\sproblem{1.21}{如题 1.21 图所示，A 为平凸透镜，B 为平玻璃板，C 为金属柱，D 为框架，A、B 间有孔隙，图中绘出的是接触的情况，而 A 固结在框架的边缘上。温度变化时，C 发生伸缩，而假设 A、B、D 都不发生伸缩。以波长 $632.8\,\mathrm{nm}$ 的激光垂直照射，试问：
\begin{enumerate}
  \item 在反射光中观察时，看到牛顿环条纹移向中央，这表示金属柱 C 的长度在增加还是减小？
  \item 若观察到有 $10$ 个亮条纹移到中央而消失，试问 C 的长度变化了多少毫米？
\end{enumerate}\fig[0.44]{1_21}}
\answerarea

\sproblem{1.22}{研究性课题：请查阅互联网，进一步了解光的干涉应用的新进展，并写成评述性论文。}
\answerarea
\clearpage
